%% kapa: exam=1 round=31 year=2020 period=2 form=A subject=appraisal_law subject_ko=감정평가관계법규 seq=17 number=17
%% source: Q-Net 감정평가사 기출문제 (cst00309) artlSeq=5207219 file=extracted/1008/5207219/A형/A형/2020년도 제31회 감정평가사 1차 시험_2교시_A형.pdf
%% conversion: kordoc 4.24.51 -> tools/gen1.py (auto); review=proofread-sonnet; flags=none
\begin{question}{17}
부동산 가격공시에 관한 법령상 비주거용 부동산가격의 공시에 관한 설명으로 옳지 않은 것은?
\begin{choices}
  \item 국토교통부장관이 비주거용 표준부동산을 선정할 경우 미리 해당 비주거용 표준부동산이 소재하는 시ㆍ도지사의 의견을 들어야 하나, 이를 시장ㆍ군수ㆍ구청장의 의견으로 대신할 수 있다.
  \item 국토교통부장관은 중앙부동산가격공시위원회의 심의를 거쳐 비주거용 표준부동산가격을 공시할 수 있다.
  \item 비주거용 표준부동산가격의 공시에는 비주거용 표준부동산의 대지면적 및 형상이 포함되어야 한다.
  \item 국토교통부장관은 비주거용 개별부동산가격의 산정을 위하여 필요하다고 인정하는 경우에는 비주거용 부동산가격비준표를 작성하여 시장ㆍ군수 또는 구청장에게 제공하여야 한다.
  \item 공시기준일이 따로 정해지지 않은 경우, 비주거용 집합부동산가격의 공시기준일은 1월 1일로 한다.
\end{choices}
\end{question}
\answer{1}
